\uuid{p3O1}
\titre{Ensembles}
\auteur{bodin ; croizat ; sacre}
\organisation{exo7}
\licence{CC BY-NC-SA 4.0}
\datecreate{2022-03-08}
\niveau{L1}
\module{Algèbre}
\chapitre{Logique, ensemble, raisonnement}
\sousChapitre{Ensemble}
\difficulte{3}
\theme{Ensembles}
\exo7qcmid{hilisit.logique.038}

\contenu{
\begin{question}
\qclassification{100.02}
Les fonctions $f$ et $g$ définies par les expressions suivantes sont-elles bijections réciproques l'une de l'autre ?
(On ne se préoccupera pas des ensembles de départ et d'arrivée.)
\begin{answers}
	\good{$f(x) = \exp(-3x)$ et $g(x) = -\frac13\ln(x)$}
	\bad{$f(x) = \cos(x+1)$ et $g(x) = \frac 1{\cos(x)}-1$}
	\good{$f(x) = \frac{x}{1+x}$ et $g(x) = \frac{x}{1-x}$}
	\bad{$f(x)= \sqrt{x+1}$ et $g(x) = x^2+1$}
\end{answers}
\begin{explanations}
Si $f(x) = \exp(-3x)$ et $g(x) = -\frac13\ln(x))$ alors
$$(f \circ g)(x)
= e^{(-3) \times (-\frac13\ln(x)) }
= e^{ \ln(x)  }
= x$$
On vérifie de même que $(g\circ f)(x)=x$, donc $f$ et $g$ sont bijections réciproques l'une de l'autre.

Si $f(x) = \cos(x+1)$ et $g(x) =  \frac 1{\cos(x)}-1$ alors $(f \circ g)(x) = \cos( \frac 1{\cos(x)})$ ne se simplifie pas du tout et n'a aucune chance d'être égal à $x$. Prendre $x=10$, alors $(f \circ g)(x)$ est une valeur renvoyée par un cosinus donc est compris entre $-1$ et $+1$ donc ne peut pas être égale à $10$.

Si $f(x) = \frac{x}{1+x}$ et $g(x) = \frac{x}{1-x}$
$$(f \circ g)(x)
= f\left(\frac{x}{1-x} \right)
=  \frac{\frac{x}{1-x}}{1+\frac{x}{1-x}}=\frac{\frac{x}{1-x}}{\frac{1-x}{1-x}+\frac x{1-x}}=\frac x{1-x+x} = x.$$
On vérifie de même que $(g \circ f)(x) = x$. Donc $f$ et $g$ sont bijections réciproques l'une de l'autre.

Si $f(x)= \sqrt{x+1}$ et $g(x) = x^2+1$ alors
$$(f \circ g)(x)
= f\left(x^2+1 \right)
= \sqrt{ x^2+2 }
$$
qui n'est pas égal à $x$ (par exemple pour $x=0$, $(f\circ g)(0)= \sqrt{2} \neq 0$).
\end{explanations}
\end{question}
}
